\subsection{投射维数、内射维数、同调维数}

设 A 为环，M 为 A-模。回忆 (A, X, 第 134 页, 定义 1)，M 的投射维数记作 \(\mathrm{dp}_{\mathrm{A}}(\mathrm{M})\)，它是 M 的诸投射分解的长度组成的集合在 \(\overline{\mathbf{Z}}\) 中的下确界。我们有 \(\mathrm{dp}_{\mathrm{A}}(0) = -\infty\)，且当 M 非零时 \(\mathrm{dp}_{\mathrm{A}}(\mathrm{M}) \geqslant 0\)。为使 M 是投射的，必须且只需 \(\mathrm{dp}_{A}(\mathrm{M}) \leqslant 0\)。

例. 设 J 为 A 的一个由 A-正则序列 \(\mathbf{x} = (x_{1}, \ldots, x_{r})\) 生成的理想。我们有 dp \(_{A}\) (A/J) ≤ r。事实上，若 A 为零环这是显然的；否则，Koszul 复形 \(\mathbf{K}_{\bullet}(\mathbf{x}, \mathrm{A})\) 是 A/J 的一个长度为 r 的自由分解（同上引文, 第 159 页, 注 3）。此外，对每个 A-模 N，A-模 Ext \(_{A}^{r}\) (A/J, N) 与 N/JN 同构（同上引文）；因此，为使 dp \(_{A}\) (A/J) = r，必须且只需 J 不等于 A（同上引文, 第 134 页, 命题 1）。

类似地，M 的内射维数记作 \(\mathrm{di}_{\mathrm{A}}(\mathrm{M})\)，定义为 M 的诸内射分解的长度组成的集合的下确界。我们有 \(\mathrm{di}_{\mathrm{A}}(0) = -\infty\)，且 \(\mathrm{di}_{\mathrm{A}}(\mathrm{M}) \geqslant 0\)（若 \(M \neq 0\)）。为使 M 是内射的，必须且只需 \(\mathrm{di}_{\mathrm{A}}(\mathrm{M}) \leqslant 0\)。

命题 1.- 设 A 为环，M 为 A-模，n 为一个 \(\geqslant 0\) 的整数。下列条件等价：

\begin{enumerate}
\def\labelenumi{(\roman{enumi})}
\item
  有 \(\mathrm{di}_{\mathrm{A}}(\mathrm{M}) \leqslant n\)（换言之，M 具有一个长度 \(\leqslant n\) 的内射分解）；
\item
  对每个 A-模 N 及每个整数 \(i > n\)，有 \(\mathrm{Ext}_A^i (\mathrm{N},\mathrm{M}) = 0\)；
\item
  对 A 的每个理想 \(\mathfrak{a}\)，有 \(\operatorname{Ext}_{\mathrm{A}}^{n+1}(\mathrm{A}/\mathfrak{a},\mathrm{M})=0\)；
\item
  对 A-模的每个正合列
\end{enumerate}

\[
0 \rightarrow \mathrm{M} \longrightarrow \mathrm{I} ^ {0} \longrightarrow \mathrm{I} ^ {1} \longrightarrow \dots \longrightarrow \mathrm{I} ^ {n - 1} \longrightarrow \mathrm{Q} \rightarrow 0,
\]

其中诸 \(I^{i}\) 是内射的，A-模 Q 是内射的。

(i)⇒(ii)：这由 A, X, 第 100 页, 定理 1 得出。

(ii)⇒(iii)：这是显然的。

\((\mathrm{iii}) \Rightarrow (\mathrm{iv}) : \text{在 (iv) 的情形中，对每个 A-模 N，存在一个同构，它把 Ext}_A^1(\mathrm{N}, \mathrm{Q}) \text{映到 Ext}_A^{n+1}(\mathrm{N}, \mathrm{M}) (\mathrm{A}, \mathrm{X}, \text{p. 128, cor. 4}) ; \text{在假设 (iii) 下，A-模 Ext}_A^1(A / \mathfrak{a}, \mathrm{Q}) \text{ 对每个理想 } \mathfrak{a} \text{ （属于 A）均为零，且 Q 是内射的 (A, X, p. 93, prop. 11).}\)

(iv)⇒(i)：考察正合列 (A, X, 第 52 页)

\[
0 \longrightarrow \mathrm{M} \longrightarrow \mathrm{I} ^ {0} (\mathrm{M}) \longrightarrow \dots \longrightarrow \mathrm{I} ^ {n - 1} (\mathrm{M}) \longrightarrow \mathrm{K} ^ {n - 1} (\mathrm{M}) \longrightarrow 0;
\]

若条件 (iv) 成立，则 A-模 \(K^{n-1}(M)\) 是内射的，由此得到 (i)。

回忆 (A, X, 第 138 页, 定义 2)，环 A 的同调维数记作 dh(A)，它是 \(\overline{\mathbf{Z}}\) 中这样的整数 \(n\) 组成的集合的上确界：存在两个 A-模 M 与 N 使 Ext \(_{A}^{n}\) (M,N) 非零。它因此也是所有 A-模的投射（或内射）维数组成的集合的上确界；当 A 是 Noether 环时，可以只限于有限生成 A-模 (A, X, 第 139 页, 推论)。

